In this section we will develop techniques to classify isomorphisms for spaces of functions with different algebraic structures. As in the preceding sections, we will be interested mostly in spaces of continuous functions between topological spaces, however the initial notions we will deal with can be defined in purely set-theoretical terms.

Let $X$ and $H_X$ be sets, and consider a class $\mathcal{A}(X)\subseteq(H_X)^X$ of $H_X$-valued functions on $X$. Given a point $x\in X$, denote by $\mathcal{A}(X)|_x$ the set of images of $x$ under elements of $\mathcal{A}(X)$, i.e.
\[\mathcal{A}(X)|_x=\left\{f(x):f\in\mathcal{A}(X)\right\}.\ntag\label{definitionimageofpointunderclassoffunction}\]

If $Y$ is another set and $\phi:Y\to X$ is a map, denote by
\[Y\times_{(\phi,\mathcal{A}(X))} H_X=\bigcup_{y\in Y}\left\{y\right\}\times\mathcal{A}(X)|_{\phi(y)}.\]

Note that $Y\times_{(\phi,\mathcal{A}(X))} H_X$ is equal to $Y\times H_X$ if and only if the following property is satisfied: For every $y\in Y$ and every $c\in H_X$, there exists $f\in\mathcal{A}(X)$ such that $f(\phi(y))=c$.

\begin{definition}\label{definitionttransform}
Let $X$, $H_X$, $Y$ and $H_Y$ be sets, $\phi:Y\to X$ be a function, and consider a class of functions $\mathcal{A}(X)\subseteq (H_X)^X$.

Given maps $\phi:Y\to X$ and $\chi:Y\times_{(\phi,\mathcal{A}(X))} H_X\to H_Y$, we define $T_{(\phi,\chi)}:\mathcal{A}(X)\to (H_Y)^Y$ by
\[(T_{(\phi,\chi)}f)(y)=\chi(y,f(\phi(y)),\qquad\forall f\in\mathcal{A}(X),\quad\forall y\in Y.\ntag\label{definitiontphichi}\]
\end{definition}

\begin{definition}
Let $X$, $H_X$, $Y$ and $H_Y$ be sets, $\phi:Y\to X$ be a map, and consider classes of functions $\mathcal{A}(X)\subseteq (H_X)^X$ and $\mathcal{A}(Y)\subseteq (H_Y)^Y$. A map $T:\mathcal{A}(X)\to\mathcal{A}(Y)$ is called \emph{$\phi$-basic} if there exists $\chi:Y\times_{(\phi,\mathcal{A}(X))} H_X\to H_Y$ such that $T=T_{(\phi,\chi)}$. We call such $\chi$ a \emph{$(\phi,T)$-transform}.
\end{definition}

\begin{remark}
In the definition above, we ignore the fact that the codomain of $T$ is $\mathcal{A}(Y)$, while the codomain of $T_{(\phi,\chi)}$ is $(H_Y)^Y$.
\end{remark}

In simpler terms, a basic map is one that is induced naturally by the transformation $\phi:Y\to X$, and the field of partial transformations $\chi(y,\cdot):\mathcal{A}(X)|_{\phi(y)}\subseteq H_X\to H_Y$ ($y\in Y$).

\begin{example}
Let $\phi:Y\to X$ and $\psi:H_X\to H_Y$ be functions. Then the map
\[T:(H_X)^X\to (H_Y)^Y,\qquad Tf=\psi\circ f\circ \phi\]
is $\phi$-basic, and the $(\phi,T)$-transform $\chi$ is given by $\chi(y,z)=\psi(z)$.
\end{example}

The next example will appear, in some form, in most applications (Sections \ref{sectionrecoveringknownresults} and \ref{sectionnewresults}).

\begin{example}
Suppose that $X$ is a locally compact Hausdorff space, $H_X$ is a Hausdorff space, $\theta_X\in C(X,H_X)$ and $\mathcal{A}(X)\subseteq C_c(X,\theta_X)$. Let $Y$ and $H_Y$ be topological spaces, $\phi:Y\to X$ be a homeomorphism, and $\chi:Y\times H_X\to H_Y$ be a continuous map such that for every $y\in Y$, the section $\chi(y,\cdot):H_X\to H_Y$ is a bijection.

For every $f\in\mathcal{A}(X)$, define $Tf\in C(Y,H_Y)$ as
\[Tf:Y\to H_Y,\qquad Tf(y)=\chi(y,f(\phi(y))),\]
and let $\mathcal{A}(Y)=\left\{Tf:f\in\mathcal{A}(X)\right\}$. Also define $\theta_Y=T\theta_X$. Then
\begin{enumerate}
\item $T$ is $\phi$-basic, and the $(\phi,T)$-transform is the restriction of $\chi$ to $Y\times_{(\phi,\mathcal{A}(X))} H_X$;
\item $(X,\theta_X,\mathcal{A}(X))$ is (weakly) regular if and only if $(Y,\theta_Y,\mathcal{A}(Y))$ is (weakly) regular. In this case, $T$ is a $\perpp$-isomorphism and $\phi$ is the $T$-homeomorphism.
\end{enumerate}
\end{example}

Note that not every $\perpp$-isomorphism is given as in the previous example.

\begin{example}
Suppose that $X=Y$ is compact Hausdorff, $H_X=H_Y=\mathbb{R}$ and $\theta_X=\theta_Y=0$, so we simply write $C(X)=C(X,\mathbb{R})$. Let $T:C(X)\to C(X)$ be any bijection satisfying $[f\neq 0]=[Tf\neq 0]$ for all $f\in C(X)$. Then $T$ is a $\perpp$-isomorphism, and the $T$-homeomorphism is the identity $\id_X:X\to X$. Let us look at two particular cases:
\begin{itemize}
\item Suppose that $X$ is not a singleton, $T(1)=2$, $T(2)=1$, and $Tf=f$ for every $f\neq 1,2$. Then $T$ is non-basic (see Proposition \ref{propositionbasicisomorphisms}(a)) and discontinuous with respect to either the topology of uniform convergence, or the topology of pointwise convergence.

\item If $X=\{\ast\}$ is a singleton, we identify $C(X)$ with $\mathbb{R}$, so any self-bijection $T:\mathbb{R}\to\mathbb{R}$ preserving $0$ is a basic $\perpp$-automorphism. In this case, the $T$-transform $\chi$ coincides with $T$ (or more precisely $\chi(\ast,z)=T(z)$ for all $z\in\mathbb{R}$), and ``most'' (cardinality-wise) of these are discontinuous: indeed, there are $\mathfrak{c}^\mathfrak{c}=2^\mathfrak{c}$ self-bijections of $\mathbb{R}\setminus\left\{0\right\}$ (where $\mathfrak{c}=2^{\aleph_0}$ is the continuum) but only $2^{\aleph_0}=\mathfrak{c}$ of these are continuous.
\end{itemize}
\end{example}

In the next proposition, we again consider only sets (without topologies).

\begin{proposition}\label{propositionbasicisomorphisms}
Let $\mathcal{A}(X)\subseteq (H_X)^X$ and $\mathcal{A}(Y)\subseteq (H_Y)^Y$, and consider maps $\phi:Y\to X$ and $T:\mathcal{A}(X)\to\mathcal{A}(Y)$. Then
\begin{enumerate}
\item[(a)] $T$ is $\phi$-basic if and only if for all $y\in Y$,
\[f(\phi(y))=g(\phi(y))\Longrightarrow Tf(y)=Tg(y),\qquad\forall f,g\in\mathcal{A}(X).\ntag\label{equationpropositionbasicisomorphisms}\]
\end{enumerate}
In this case,
\begin{enumerate}
\item[(b)] the $(\phi,T)$-transform $\chi$ is unique.
\item[(c)] A section $\chi(y,\cdot)$ is injective if and only if
\[Tf(y)=Tg(y)\Longrightarrow f(\phi(y))=g(\phi(y)),\qquad\forall f,g\in\mathcal{A}(X).\ntag\label{equationpropositionbasicisomorphisms2}\]
\item[(d)] A section $\chi(y,\cdot)$ is surjective if and only if
\[H_Y=\left\{Tf(y):f\in\mathcal{A}(X)\right\}.\]
\end{enumerate}
\end{proposition}
\begin{proof}
If $T$ satisfies (\ref{equationpropositionbasicisomorphisms}), define $\chi$ by
\[\chi(y,t)=Tf(y),\]
whenever $f\in\mathcal{A}(X)$ is any function satisfying $f(\phi(y))=t$. Then $\chi(y,t)$ does not depend on the choice of $f$ by Equation (\ref{equationpropositionbasicisomorphisms}), and hence $T$ is $\phi$-basic.

The converse direction, as well as items (b), (c) and (d) are immediate from the formula $Tf(y)=\chi(y,f(\phi(y)))$, which holds for all $f\in\mathcal{A}(X)$ and $y\in Y$.\qedhere
\end{proof}

The next proposition proves that, under usual regularity hypotheses, $\perpp$-\hspace*{0pt}isomorphisms can be basic only with respect to their corresponding homeomorphisms.

\begin{proposition}\label{propositionuniquemapwhichmakesbasicisthomeomorphism}
Let $X$ and $Y$ be locally compact Hausdorff spaces, $H_X$ and $H_Y$ be Hausdorff spaces, $\theta_X\in C(X,H_X)$, $\theta_Y\in C(Y,H_Y)$ and $\mathcal{A}(X)\subseteq C_c(X,\theta_X)$ and $\mathcal{A}(Y)\subseteq C_c(Y,\theta_Y)$ be two regular classes of functions.

Suppose that $T:\mathcal{A}(X)\to\mathcal{A}(Y)$ is a $\perpp$-isomorphism which is $\phi$-basic for some map $\phi:Y\to X$, and such that $T^{-1}:\mathcal{A}(Y)\to\mathcal{A}(X)$ is $\psi$-basic for some map $\psi:X\to Y$. Then $\phi$ is invertible, $\psi=\phi^{-1}$ and $\phi$ is the $T$-homeomorphism.
\end{proposition}

\begin{lemma}\label{lemmapropositionuniquemapwhichmakesbasicisthomeomorphism}
Under the hypotheses of Proposition \ref{propositionuniquemapwhichmakesbasicisthomeomorphism}, $\phi:Y\to X$ is the only map for which $T$ is $\phi$-basic.
\end{lemma}
\begin{proof}
Let $\phi_1=\phi$ and assume that $\phi_2:Y\to X$ is another map such that $T$ is $\phi_2$-basic. Given $y\in Y$, we will prove that $\phi_1(y)=\phi_2(y)$. Assume that this is not true, and let us deduce a contradiction. Note, in particular, that $|Y|\geq 2$, so regularity of $\mathcal{A}(Y)$ implies $|H_Y|\geq 2$. Let $\chi_1$ and $\chi_2$ be the $(\phi_1,T)$ and $(\phi_2,T)$-transforms, respectively.

Since $\mathcal{A}(Y)$ is regular and $|H_Y|\geq 2$, there exists $c\in H_X$ such that
\[\chi_1(y,c)\neq\theta_Y(y).\ntag\label{inequalitylemmapropositionuniquemapwhichmakesbasicisthomeomorphism}\]
Since we are assuming that $\phi_1(y)\neq\phi_2(y)$, regularity of $\mathcal{A}(X)$ yields $f\in\mathcal{A}(X)$ such that $f(\phi_1(y))=c$ and $f(\phi_2(y))=\theta_X(\phi_2(y))$. As $T$ is a $\perpp$-isomorphism we have $T\theta_X=\theta_Y$, and the properties of the $(\phi_i,T)$-transforms $\chi_i$ ($i=1,2$) imply
\begin{align*}
\theta_Y(y)&=T\theta_X(y)=\chi_2(y,\theta_X(\phi_2(y)))=\chi_2(y,f(\phi_2(y)))\\
&=Tf(y)=\chi_1(y,f(\phi_1(y)))=\chi_1(y,c)
\end{align*}
contradicting inequality (\ref{inequalitylemmapropositionuniquemapwhichmakesbasicisthomeomorphism}).
\end{proof}

\begin{proof}[Proof of Proposition \ref{propositionuniquemapwhichmakesbasicisthomeomorphism}]
Since $T$ and $T^{-1}$ are $\phi$ and $\psi$-basic, respectively, then the identity map $\id_{\mathcal{A}(X)}=T^{-1}\circ T$ is $(\psi\circ\phi)$-basic (by Proposition \ref{propositionbasicisomorphisms}(a)). However, $\id_{\mathcal{A}(X)}$ is $\id_Y$-basic, (again by Proposition \ref{propositionbasicisomorphisms}(a); the $(\id_Y,\id_{\mathcal{A}(X)})$-transform is given by $\chi(y,t)=t$). Lemma \ref{lemmapropositionuniquemapwhichmakesbasicisthomeomorphism} applied to $T^{-1}\circ T$ then implies $\psi\circ\phi=\id_Y$, and similarly $\phi\circ\psi=\id_X$. This proves that $\phi^{-1}=\psi$.

Now denote by $\rho:Y\to X$ the $T$-homeomorphism, so we need to prove that $\phi=\rho$. Proposition \ref{propositionbasicisomorphisms}(a) applied to $T$ and $T^{-1}$ yields, for all $y\in Y$ and all $f\in\mathcal{A}(X)$,
\[Tf(y)=\theta_Y(y)\iff f(\phi(y))=\theta_X(\phi(y)).\ntag\label{inequalitylemmapropositionuniquemapwhichmakesbasicisthomeomorphism2}\]
Given $y\in Y$, regularity of $\mathcal{A}(X)$ gives us
\[\left\{\phi(y)\right\}=\bigcap_{f(\phi(y))\neq\theta_X(\phi(y))}[f\neq\theta_X]\]
and using the equivalence (\ref{inequalitylemmapropositionuniquemapwhichmakesbasicisthomeomorphism2}), the definition of $\sigma^{\theta_X}$ (Definition \ref{definitionsigma}) and the characterization of the $T$-homeomorphism $\rho$ (Definition \ref{definitionthomeomorphism}), we obtain
\[\left\{\phi(y)\right\}\subseteq\bigcap_{Tf(y)\neq\theta_Y(y)}\sigma^{\theta_X}(f)=\bigcap_{Tf(y)\neq\theta_Y(y)}\rho(\sigma^{\theta_Y}(Tf))=\left\{\rho(y)\right\}\]
where the last equality follows as $\rho$ is injective and $\mathcal{A}(Y)=T(\mathcal{A}(X))$ is regular. We thus conclude that $\phi=\rho$ is the $T$-homeomorphism.\qedhere
\end{proof}

\subsection{Algebraic signatures and basic maps}

In our examples and consequences we will consider different algebraic structures on spaces of continuous functions. To this end, recall (see \cite{MR1221741}) that an \emph{algebraic signature}\index{Signature!Algebraic} is a collection $\eta$ of pairs $(*,n)$, where $*$  is a (function) symbol and $n$ is a non-negative integer, called the \emph{arity} of $*$. A \emph{model} of $\eta$ consists of a set $H$ and a map associating to each $(*,n)\in\eta$ a function $*:H^n\to H$, $(c_1,\ldots,c_n)\mapsto c_1*\cdots*c_n$. (We use the convention that $H^0$ is a singleton set, so that a $0$-ary function symbol is the same as a constant.)

This can be seen as a particular case of a metric signature (Definition \ref{definitionmetricsignature}), where only discrete metric spaces (with the $0-1$ metric) are considered. (See Example \ref{examplesofmetricsignatures}(2).)

If $H$ is a model of $\eta$ and $X$ is a set then the function space $H^X$ can also be regarded as a model of $\eta$ with the pointwise structure: $(f_1*\cdots*f_n)(x)=f_1(x)*\cdots*f_n(x)$ for all $f_1,\ldots,f_n\in H^X$, all $x\in X$, and all $n$-ary function symbols $\ast$.

A \emph{morphism} of two models $H_1$ and $H_2$ of a given signature $\eta$ is a map $m:H_1\to H_2$ such that for any $n$-ary function symbol $*$ of $\eta$ and any $x_1,\ldots,x_n\in H_1$, we have $m(x_1*\cdots*x_n)=m(x_1)*\cdots*m(x_n)$.

Finally, a \emph{submodel} of a model $H$ of a signature $\eta$ is a subset $K\subseteq H$ such that for all $n$-ary symbols $*$ of $\eta$ and any $d_1,\ldots,d_n\in K$, $d_1*\cdots*d_n\in K$, so that $K$ can be naturally regarded as a model of $\eta$.

In the topological setting, a \emph{continuous model} $H$ of a signature $\eta$ is defined in the same manner, but we assume that all maps are continuous. In this case, if $X$ is a topological space then $C(X,H)$ is a submodel of $H^X$.

\begin{example}
The standard signature for groups consists of a binary ($2$-ary) operation $\cdot$ (the group operation), a unary operation $(\ )^{-1}$ (the inversion), and a constant ($0$-ary operation) $1$ (the unit).
\end{example}

\begin{proposition}\label{propositionmodelmorphism}
Let $X$ and $Y$ be sets. Suppose that $H_X$ and $H_Y$ are models for an algebraic signature $\eta$, and that $\mathcal{A}(X)$ and $\mathcal{A}(Y)$ are submodels of $(H_X)^X$ and $(H_Y)^Y$. Then for all $x\in X$, $\mathcal{A}(X)|_x$ is a submodel of $H_X$, and similarly $\mathcal{A}(Y)|_y$ is a submodel of $H_Y$ for all $y\in Y$. (See Equation (\ref{definitionimageofpointunderclassoffunction}).)

Let $T:\mathcal{A}(X)\to\mathcal{A}(Y)$ be a basic map with respect to a function $\phi:Y\to X$, and let $\chi$ be the $T$-transform. Then $T$ is a morphism (for $\eta$) if and only if every section $\chi(y,\cdot)$ is a morphism (from $\mathcal{A}(X)|_{\phi(y)}$ to $\mathcal{A}(Y)|_y$).
\end{proposition}
\begin{proof}
Given $x\in X$, the evaluation map $\pi_x:(H_X)^X\to H_X$, $\pi_x(f)=f(x)$, is a morphism, and it follows that $\mathcal{A}(X)|_x=\pi_x(\mathcal{A}(X))$ is a submodel of $H_X$.

Suppose that $T$ as above is a morphism and let $y\in Y$. Given an $n$-ary symbol $\ast$ of $\eta$, suppose $t_1,\ldots,t_n\in \mathcal{A}(X)_{\phi(y)}$. Take functions $f_1,\ldots,f_n\in\mathcal{A}(X)$ such that $f_i(\phi(y))=t_i$, so $f_1\ast\cdots\ast f_n\in\mathcal{A}(X)$ since $\mathcal{A}(X)$ is a submodel of $C(X,H_X)$. Then $(f_1*\cdots*f_n)(y)=t_1*\cdots*t_n$, and thus
\begin{align*}
\chi(y,t_1*\cdots*t_n)&=T(f_1*\cdots*f_n)(y)=Tf_1(y)*\cdots*Tf_n(y)\\
&=\chi(y,t_1)*\cdots*\chi(y,t_n).
\end{align*}
This proves that $\chi(y,\cdot)$ is a morphism. The converse implication is similar.\qedhere
\end{proof}

%In the cases of interest, we will consider $\perpp$-isomorphisms $T$ which are basic with respect to the $T$-homeomorphism.

\subsection{Group-valued maps}\label{subsectiongroupvaluedmaps}

We first note a simple property for basic isomorphisms of groups of functions which will be used throughout the next section:

\begin{proposition}\label{propositionbasicnessofgroupvalued}
Suppose that $H_X$ and $H_Y$ are groups, $\mathcal{A}(X)$ and $\mathcal{A}(Y)$ are subgroups of $(H_X)^X$ and $(H_Y)^Y$, respectively, $T:\mathcal{A}(X)\to\mathcal{A}(Y)$ is a group isomorphism and $\phi:Y\to X$ is a function. Then $T$ is basic with respect to $\phi$ if and only if for all $y\in Y$,
\[f(\phi(y))=1\Longrightarrow Tf(y)=1,\qquad\forall f\in\mathcal{A}(X).\]
\end{proposition}
\begin{proof}
This follows directly from \ref{propositionbasicisomorphisms}(a), that $f(x)=g(x)$ if and only if $(fg^{-1})(x)=1$ and the fact that $T$ preserves products and inverses.\qedhere
\end{proof}

In the topological setting, assume that $H$ is a topological group, $X$ is a locally compact Hausdorff space and $\theta\in C(X,H)$. Then, $C_c(X,\theta)$ and $C_c(X,1)$ are $\perpp$-isomorphic, via $f\mapsto f\theta^{-1}$, and in fact this same map can be used to show that $C_c(X,\theta)$ is (weakly) regular if and only if $C_c(X,1)$ is (weakly) regular.

Moreover, if $C_c(X,\theta)$ is a subgroup of $C(X,H)$ then $\theta=1$ outside of a compact set. Indeed, let $f,g\in C_c(X,\theta)$ be arbitrary, and suppose that $fg\in C_c(X,\theta)$ as well. As $f$ and $fg$ coincide with $\theta$ outside of a compact, then $g=1$ outside of a compact. Since $g=\theta$ outside of a compact, then $\theta=1$ outside of a compact. Therefore, in this case, we obtain $C_c(X,\theta)=C_c(X,1)$.

Since we will be interested in structure-preserving maps between submodels of $C(X,H)$, in the case the codomain is a group we may always assume that $\theta=1$. In the case that $H=\mathbb{R}$ or $\mathbb{C}$, as additive groups, we reobtain the usual notion of support.

\subsection{Continuity}

Now, we study continuity of basic $\perpp$-isomorphisms and relate it to the continuity of its transform, and for this we need a few results from general topology.

\begin{proposition}\label{propositiondisjointopensets}
If $F$ is an infinite subset of a regular Hausdorff space $X$, then there exists a countable infinite subset $\left\{y_1,y_2,\ldots\right\}\subseteq F$ and open sets $U_n$ such that $y_n\in U_n$ for all $n$ and $U_n\cap U_m=\varnothing$ if $n\neq m$.
\end{proposition}
\begin{proof}
First assume $F$ has a cluster point $y$. Let $y_1\in Y\setminus\left\{y\right\}$ be arbitrary. Take disjoint neighbourhoods $U_1$ of $y_1$ and $V_1$ of $y$, and then repeat the procedure in $V_1$. This way, we can construct sequences of points $y_1,y_2,\ldots\in Y$ and of neighbourhoods $U_n$ of $y_n$ and $V_n$ of $y$ such that
\begin{itemize}
\item $V_{n+1}\subseteq V_n$;
\item $U_n\cap V_n=\varnothing$;
\item $U_{n+1}\subseteq V_n$.
\end{itemize}
If $n<m$, then
\[U_m\cap U_n\subseteq V_{m-1}\cap U_n\subseteq V_n\cap U_n=\varnothing,\]
as wanted.

Now assume that $F$ does not have a cluster point. Take any infinite countable subset $\left\{y_n:n\in\mathbb{N}\right\}\subseteq F$. Given $k$, since $y_k$ is a not a cluster point of $F$, we can take neighbourhoods $V_k$ such that $V_k\cap\left\{y_n:n\in\mathbb{N}\right\}=\left\{y_k\right\}$. By regularity, take a neighbourhood $U_1$ of $y_1$ such that $\overline{U_1}\subseteq V_1$, and then inductively take, for a given $n$, a neighbourhood $U_{n+1}$ of $y_{n+1}$ such that $\overline{U}_{n+1}\subseteq V_{n+1}\setminus\bigcup_{i=1}^n\overline{U_i}$.\qedhere
\end{proof}

For the next proposition, recall (\cite[27.4]{MR0264581}) that a topological space $H$ is \emph{locally path-connected} if every point $t\in H$ admits a neighbourhood basis consisting of path-connected subsets.

\begin{proposition}\label{propositionconstructfunction}
Let $X$ be a locally compact Hausdorff space, $\left\{x_n\right\}_n$ be a sequence of elements of $X$, $\left\{U_n\right\}_n$ a sequence of open subsets of $X$ such that
\begin{enumerate}
\item[(i)] $x_n\in U_n$;
\item[(ii)] $U_n\cap U_m=\varnothing$ for all $n\neq m$.
\end{enumerate}
Let $H$ be a Hausdorff first-countable locally path-connected topological space and consider $\left\{g_n:U_n\to H\right\}$ a family of continuous functions such that $g_n(x_n)$ converges to some $t\in H$. Then
\begin{enumerate}
\item[(a)] there exists a continuous function $f:X\to H$ such that $f(x_n)=g_n(x_n)$ for all sufficiently large $n$, and $f(x)=t$ for all $x\not\in\bigcup_n U_n$.
\item[(b)] if $H=\mathbb{R}$, there is a continuous function $f:X\to H$ such that $f=g_n$ on a neighbourhood of $x_n$ and $f(x)=t$ for all $x\not\in\bigcup_n U_n$.
\end{enumerate}
(Note that it is not necessary to take subsequences!)
\end{proposition}
\begin{proof}
\begin{enumerate}
\item[(a)] Let $\left\{W_n\right\}_n$ be a decreasing basis of path-connected neighbourhoods of $t$. Disregarding any $n$ such that $g_n(x_n)$ does not belong to $W_1$, and repeating the sets $W_k$ if necessary (i.e., considering a new sequence of neighbourhoods of $t$ of the form
\[W_1,W_1,\ldots,W_1,W_2,W_2,\ldots,W_2,\ldots,\]
where each $W_k$ is repeated finitely many times) we may assume that $t_n:=g_n(x_n)\in W_n$.

For each $n$, take a continuous path $\alpha_n:[0,1]\to W_n$ such that $\alpha_n(0)=t_n$ and $\alpha_n(1)=t$. Now take continuous functions $b_n:X\to[0,1]$ such that $b_n(x_n)=0$ and $b_n=1$ outside $U_n$. Define $f$ as $\alpha_n\circ b_n$ on each $U_n$, and as $t$ on $X\setminus\bigcup_n U_n$.

Let us show that $f$ is continuous: If $x\in U_n$ for some $n$ then $f=\alpha_n\circ b_n$ on a neighbourhood of $x$ and so it is continuous at $x$. Let us then assume that $x\not\in\bigcup_n U_n$ and in particular $\alpha_n(b_n(x))=f(x)=t$ for all $n$.

Consider a neighbourhood $W_N$ of $t$, where $N$ is an integer. Take a neighbourhood $U$ of $x$ such that if $n\leq N$ and $y\in U$, then $\alpha_n(b_n(y))\in W_N$ (e.g. $U=\bigcap_{n=1}^N(\alpha_n\circ b_n)^{-1}(W_N)$). If $y\in U$, there are three cases:
\begin{itemize}
\item If $y\in U_n$ for some $n\leq N$, then $f(y)=\alpha_n(b_n(y))\in W_N$, by the choice of $U$;
\item If $y\in U_n$ for some $n>N$, then $f(y)=\alpha_n(b_n(y))\in W_n\subseteq W_N$, by the choice of $\alpha_n$;
\item If $y\not\in\bigcup_n U_n$ then $f(y)=t\in W_N$.
\end{itemize}
In any case, we have $f(U)\subseteq W_N$, which proves the continuity of $f$.
\item[(b)] For each $n$, choose an open neighbourhood $V_n$ of $x_n$ such that $\overline{V_n}\subseteq U_n$. By considering even smaller neighbourhoods if necessary we can assume $|g_n(x)-g_n(x_n)|<1/n$ for all $x\in V_n$. Up to modifying $g_n$ except on a neighbourhood of $x_n$, we can assume $g_n=t$ on $\overline{U_n}\setminus V_n$. Define $f=g_n$ on each $U_n$ and $f=t$  on $X\setminus\bigcup_n U_n$. The proof that $f$ is continuous is similar to that of item (a).\qedhere
\end{enumerate}
\end{proof}

%\begin{proposition}\label{propositioncontinuity}
%Let $X$ and $H$ be topological spaces with $H$ first countable, and $f:X\to H$ be a function. Then $f$ is continuous at $a\in X$ if and only if for every directed set $(I,\leq)$ and every converging net $y_i\to a$ ($i\in I$), there exists an increasing sequence $\left\{i(n)\right\}_{n\in\mathbb{N}}$ in $I$ such that $f(y_{i(n)})\to f(a)$.
%\end{proposition}
%\begin{proof}
%Let $\left\{U_n:n\in\mathbb{N}\right\}$ be a decreasing countable neighbourhood basis of $f(a)$.
%
%Suppose that $f$ is continuous at $a$ and let $(I,\leq)$ and $y_i$ be as above. We can inductively construct an increasing sequence $i(n)$, $n\in\mathbb{N}$, of elements of $I$ such that $f(y_j)\in U_n$ whenever $j\geq i(n)$, and this sequence satisfies the desired properties.
%
%If $f$ is not continuous at $a$, then there exists a directed set $(J,\leq)$ and a converging net $x_j\to a$ such that $f(x_j)\not\to f(a)$. Then we can find a neighbourhood $U$ of $f(a)$ and a cofinal subset $I$ of $J$ such that $f(x_i)\not\in U$  for all $i\in I$. Of course no sequence $i(n)$ as above exists in this case.\qedhere
%\end{proof}

\begin{theorem}\label{theorembasicisomorphisms}
Let $X$ and $Y$ be locally compact Hausdorff and for $Z\in\left\{X,Y\right\}$, $H_Z$ a Hausdorff space and $\theta_Z\in C(Z,H_Z)$ be given such that $(Z,\theta_Z,C_c(Z,\theta_Z))$ is regular.

Suppose that $T:C_c(X,\theta_X)\to C_c(Y,\theta_Y)$ is a $\perpp$-isomorphism, that $\phi:Y\to X$ is the $T$-homeomorphism $\phi$, and that that $T$ is $\phi$-basic. Let $\chi:Y\times H_X\to H_Y$ be the corresponding $(\phi,T)$-transform. Consider the following statements:
\begin{enumerate}
\item[(1)] $\chi$ is continuous.
\item[(2)] Each section $\chi(y,\cdot)$ is a continuous;
\item[(3)] $T$ is continuous with respect to the topologies of pointwise convergence.
\end{enumerate}
Then the implications (1)$\Longrightarrow$(2)$\iff$(3) always hold.

If $X$, $Y$ and $H_X$ are first countable, $H_X$ is locally path-connected and $\theta_X$ is constant, then $(2)\Longrightarrow (1)$.
\end{theorem}

\begin{remark}
\begin{enumerate}
\item In the last part of the theorem, if $H_X$ admits any structure of topological group then the condition that $\theta_X$ is constant can be dropped, by the discussion in Subsection \ref{subsectiongroupvaluedmaps}.
\item The domain of the $(\phi,T)$-transform $\chi$ is $Y\times H_X$ because we assume that $C_c(X,\theta_X)$ is regular.
\end{enumerate}
\end{remark}

\begin{proof}
The implication (1)$\Rightarrow$(2) is trivial.

\begin{description}
\item[\normalfont{(2)$\Rightarrow$(3):}] Suppose $f_i\to f$ pointwise. Then for all $y$, the section $\chi(y,\cdot)$ is continuous, thus
\[Tf_i(y)=\chi(y,f_i(\phi(y)))\to\chi(y,f(\phi(y)))=Tf(y).\]
This proves that $Tf_i\to Tf$ pointwise.
	\item[\normalfont{(3)$\Rightarrow$(2):}]Assume that $T$ is continuous with respect to pointwise convergence. Let $y\in Y$ be fixed. Suppose that $t_i\to t$ in $H_X$, and let us prove that $\chi(y,t_i)\to\chi(y,t)$. Choose any function $f\in C_c(X,\theta_X)$ such that $f(\phi(y))=t$.
	
	Let $\operatorname{Fin}(X)$ be the collection of finite subsets of $X$, ordered by inclusion. We will construct a net $\left\{f_{(F,i)}\right\}_{(F,i)\in\operatorname{Fin}(X)\times I}$ of functions in $C_c(X,\theta_X)$ such that
	\begin{itemize}
		\item[(i)] $f_{(F,i)}\to f$ pointwise;
		\item[(ii)] $f_{(F,i)}(\phi(y))=t_i$.
	\end{itemize}
	Given $F\in\operatorname{Fin}(X)$ and $i\in I$, consider a family $\left\{U_x:x\in F\cup\left\{\phi(y)\right\}\right\}$ of pairwise disjoint open sets such that $x\in U_x$ for each $x$.
	
	Given $x\in F\cup\left\{\phi(y)\right\}$ consider, by regularity of $C_c(X,\theta_X)$, a function $f_{(F,i,x)}\in C_c(X,\theta_X)$ such that $\supp(f_{(F,i,x)})\subseteq U_x$ and
	\[f_{(F,i,x)}(x)=\begin{cases}t_i,&\text{if }x=\phi(y)\\
	f(x),&\text{otherwise}.\end{cases}\]
	We then define $f_{(F,i)}:X\to H_X$ by
	\[f_{(F,i)}=f_{(F,i,x)}\text{ on each set }U_x\text{ and }f_{(F,i)}=\theta_X\text{ on }X\setminus\bigcup_{x\in F\cup\left\{\phi(y)\right\}}U_x.\]
	This way we obtain $f_{(F,i)}\in C_c(X,\theta_X)$. Properties (i) and (ii) above are immediate because $t_i\to t$.
	
	Since $f_{(F,i)}\to f$ pointwise then $Tf_{(F,i)}\to Tf$ pointwise. For each $F\in\operatorname{Fin}(X)$ and $i\in I$, we have
\[\chi(y,t_i)=\chi(y,f_{(F,i)}(\phi(y)))=Tf_{(F,i)}(y)\]
so by considering $i$ and $F$ sufficiently large we see that $\chi(y,t_i)\to Tf(y)=\chi(y,t)$ as $i\to\infty$.
\end{description}

We now assume further that $X$, $Y$ and $H_X$ are first countable, $H_X$ is locally path-connected and $\theta_X$ is constant. Let $c\in H_X$ such that $\theta_X(x)=c$ for all $x\in X$.
\begin{description}
\item[\normalfont{(2)$\Rightarrow$(1):}] Assume that each section $\chi(y,\cdot)$ is continuous. In order to prove that $\chi$ is continuous, we simply need to prove that for any converging sequence $(y_n,t_n)\to (y,t)$ in $Y\times H_X$, we can take a subsequence $(y_{n'},t_{n'})$ such that $\chi(y_{n'},t_{n'})\to \chi(y,t)$ as $n'\to\infty$.

Given a converging sequence $(y_n,t_n)\to(y,t)$, consider an open $Y'\subseteq Y$ with compact closure such that $y,y_i\in Y'$ for all $i$.

We have two cases: If for a given $z\in Y$ the set $N(z)=\left\{n\in\mathbb{N}:y_n=z\right\}$ is infinite, then we necessarily have $z=y$. Restricting the sequence $(y_n,t_n)$ to $N(y)$ and using continuity of the section $\chi(y,\cdot)$, we obtain $\chi(y_n,t_n)=\chi(y,t_n)\to \chi(y,t)$ as $n\to\infty$, $n\in N(y)$.

Now assume that none of the sets $N(z)=\left\{n\in\mathbb{N}: y_n=z\right\}$ ($z\in Y$) is infinite. We may then take a subsequence and assume that all the elements $y_n$ are distinct, and actually never equal to $y$. Using Proposition \ref{propositiondisjointopensets} and taking another subsequence if necessary, consider pairwise disjoint open subsets $U_n\subseteq\phi(Y')$ with $\phi(y_n)\in U_n$ and $\phi(y)\in X\setminus\bigcup_n U_n$. We then consider, by Proposition \ref{propositionconstructfunction}, a continuous function $f:\overline{\phi(Y')}\to H_X$ such that $f(\phi(y_n))=t_n$ and $f=t$ on $X\setminus \bigcup_n U_n$. In particular, $f=t$ on the boundary $\partial(\phi(Y'))$.

We now need to extend $f$ to an element of $C_c(X,\theta_X)$ (this is where we use that $\theta_X=c$ is constant). We have two cases:%Here is where we use that $\theta_X$ is constant.
\begin{description}
\item[Case 1:] $t$ is in the path-connected component of $c$:

Since $H_X$ is locally path-connected, there is a continuous path $\beta:[0,1]\to H_X$ with $\beta(0)=t$ and $\beta(1)=c$. Let $g:X\to[0,1]$ be a function with $g=0$ on $\phi(Y')$ and $g=1$ outside of a compact containing $\phi(Y')$. By defining $f=\beta\circ g$ outside of $\phi(Y')$, we obtain $f\in C_c(X,\theta_X)$. ($f$ is continuous because $f=t=\beta\circ g$ on $\partial(\phi(Y')$.)
\item[Case 2:] $t$ is not in the path-connected component of $c$:

Since $H_X$ is locally path-connected, its path-connected components are clopen, and regularity of $C_c(X,\theta_X)$ then implies that $X$ (and thus also $Y=\phi(X)$) is zero-dimensional. In particular, we could have assumed at the beginning that $Y'$ is clopen, so simply set $f=c$ outside of $\phi(Y')$.
\end{description}

In any case, we obtain $f\in C_c(X,\theta_X)$ with $f(\phi(y))=t$ and $f(\phi(y_n))=t_n$, so
\[\chi(y_n,t_n)=Tf(y_n)\to Tf(y)=\chi(y,t).\qedhere\]
\end{description}
\end{proof}

\subsection{Non-vanishing bijections}

Let $X$ and $Y$ be compact Hausdorff spaces, $H_X$ and $H_Y$ Hausdorff spaces, $\theta_X\in C(X,H_X)$, $\theta_Y\in C(Y,\theta_Y)$ and $\mathcal{A}(X)$ and $\mathcal{A}(Y)$ regular subsets of $C_c(X,\theta_X)$ and $C_c(Y,\theta_Y)$, respectively.

\begin{definition}[\cite{MR2324919}]\label{definitionnonvanishingbijection}\index{Non-vanishing}
We call a bijection $T:\mathcal{A}(X)\to\mathcal{A}(Y)$ \emph{non-vanishing} if for every $f_1,\ldots,f_n\in\mathcal{A}(X)$,
\[\bigcap_{i=1}^n[f_i=\theta_X]=\varnothing\iff\bigcap_{i=1}^n[Tf_i=\theta_Y]=\varnothing.\]
\end{definition}

\begin{proposition}\label{propositionnonvanishingbijectionisperppisomorphism}
If $T:\mathcal{A}(X)\to\mathcal{A}(Y)$ is a non-vanishing bijection, then $T$ is a $\perpp$-isomorphism.
\end{proposition}
\begin{proof}
First note that
\begin{align*}
f\perp g&\iff[f\neq\theta_X]\cap[g\neq\theta_X]=\varnothing\iff[f=\theta_X]\cup[g=\theta_X]=X\\
&\iff\text{ every closed subset of }X\text{ intersects }[f=\theta_X]\text{ or }[g=\theta_X].
\end{align*}
So from this and the fact that the sets $[h=\theta_X]$ ($h\in\mathcal{A}(X)$) form a closed basis (and Cantor's Intersection Theorem), we obtain
\begin{align*}
f\perp g\iff&\forall h_1,\ldots,h_n\in\mathcal{A}(X),\\
&\text{if }\bigcap_{i=1}^n[h_i=\theta_X]\cap[f=\theta_X]=\bigcap_{i=1}^n[h_i=\theta_X]\cap[g=\theta_X]=\varnothing\\
&\text{then }\bigcap_{i=1}^n[h_i=\theta_X]=\varnothing.
\end{align*}
The last condition is preserved under non-vanishing bijections, and so $T$ is a $\perp$-isomorphism.

Now we show that $f\perpp g$ is equivalent to the following statement: There are finite families $\left\{a_i\right\}$, $\left\{b_j\right\}$ and $\left\{c_k\right\}$ in $\mathcal{A}(X)$ such that
\begin{enumerate}
\item[(i)] $\bigcap_{i,j,k}[a_i=\theta_X]\cap[b_j=\theta_X]\cap[c_k=\theta_X]=\varnothing$;
\item[(ii)] $a_i\perp b_j$ for all $i$ and $j$;
\item[(iii)] $f\perp b_j$, $f\perp c_k$, $g\perp a_i$, and $g\perp c_k$ for all $i$, $j$ and $k$.
\end{enumerate}
Indeed, if $f\perpp g$, by regularity and compactness one can take finite families $\left\{a_i\right\}$ and $\left\{b_j\right\}$ satisfying (ii) and such that $\supp(f)\subseteq\bigcup_i[a_i\neq\theta_X]$ and $\supp(g)\subseteq\bigcup_j[b_j\neq\theta_X]$. Then take a finite family $\left\{c_k\right\}$ such that (iii) is satisfied and such that
\[X\setminus\left(\bigcup_{i,j}[a_i\neq\theta_X]\cup[b_j\neq\theta_X]\right)\subseteq\bigcup_k[c_k\neq\theta_X],\]
which implies (i).

Conversely, if such families $\left\{a_i\right\},\left\{b_j\right\},\left\{c_k\right\}$ in $\mathcal{A}(X)$ exist, suppose $x\in\supp(f)\cap\supp(g)$. By item (i), there is at least one of the indices $i$, $j$ or $k$ such that $a_i(x)$, $b_j(x)$ or $c_k(x)$ is not equal to $\theta_X(x)$. Since $x\in\supp(f)$, and $f\perp b_j$ and $f\perp c_k$ by (iii), the only possibility is that $a_i(x)\neq\theta_X(x)$ for some $i$. The same argument with $g$ in place of $f$ yields $b_j(x)\neq\theta_X(x)$ for some $j$, contradicting (ii).

Therefore $T$ is a $\perpp$-isomorphism.\qedhere
\end{proof}

\begin{theorem}\label{theoremnonvanishing}
For every non-vanishing bijection $T:\mathcal{A}(X)\to\mathcal{A}(Y)$ there is a unique homeomorphism $\phi:Y\to X$ such that $[f=\theta_X]=\phi([Tf=\theta_Y])$ for all $f\in\mathcal{A}(X)$.
\end{theorem}
\begin{proof}
By Proposition \ref{propositionnonvanishingbijectionisperppisomorphism}, we already know that $T$ is a $\perpp$-isomorphism, so let $\phi$ be the $T$-homeomorphism. Recall (Definition \ref{definitionsigma}) that $\sigma^{\theta_X}(f)=\operatorname{int}(\overline{[f\neq\theta_X]})$ and $Z^{\theta_X}(f)=\operatorname{int}([f=\theta_X])$ for all $f\in\mathcal{A}(X)$. Let us prove that
\begin{align*}
f(x)=\theta_X(x)\iff&\forall h_1\ldots,h_n\in\mathcal{A}(X),\\
&\text{if }x\not\in\bigcup_{i=1}^n\sigma^{\theta_X}(h_i)\text{ then }\bigcap_{i=1}^n[h_i=\theta_X]\cap[f=\theta_X]\neq\varnothing,\ntag\label{equationnonvanishingproperty}
\end{align*}
Indeed, for the ``$\Rightarrow$'' direction, assume that $f(x)=\theta_X(x)$ and $h_1,\ldots,h_n$ are such that $x\not\in\bigcup_i\sigma^{\theta_X}(h_i)$. Then $x\in\bigcap_i[h_i=\theta_X]\cap[f=\theta_X]$, and this set is nonempty.

For the converse we prove the contrapositive. Assume that $f(x)\neq\theta_X(x)$, so regularity and compactness give us $h_1,\ldots,h_n\in\mathcal{A}(X)$ such that $x\in Z^{\theta_X}(h_i)$ for all $i$ and $X\setminus[f\neq\theta_X]\subseteq\bigcup_i[h_i\neq\theta_X]$, which negates the right-hand side of (\ref{equationnonvanishingproperty}).

Since $\phi(\sigma^{\theta_Y}(Th))=\sigma^{\theta_X}(h)$ for all $h\in\mathcal{A}(X)$ and $T$ is non-vanishing, the condition of  (\ref{equationnonvanishingproperty}) is preserved by $T$ and therefore, $\phi$ has the desired property.\qedhere
\end{proof}